\documentclass[a4paper,11pt]{article}
\usepackage[utf8]{inputenc}
\usepackage[T1]{fontenc}
\usepackage[ngerman]{babel}
\usepackage{amsmath,amssymb}
\usepackage{xcolor}
\usepackage{colortbl}
\usepackage{array}
\usepackage{tikz}
\usepackage{enumitem}
\usepackage{geometry}
\geometry{a4paper,margin=2cm}
\definecolor{bgdark}{HTML}{1E1E1E}
\definecolor{fgtext}{HTML}{E6E6E6}
\definecolor{accent}{HTML}{6CB6FF}
\definecolor{accent2}{HTML}{F2A65A}
\definecolor{gridcol}{HTML}{4A4A4A}
\definecolor{linecol}{HTML}{8A8A8A}
\pagecolor{bgdark}
\color{fgtext}
\arrayrulecolor{linecol}
\newcommand{\karo}[1]{\par\noindent\begin{tikzpicture}\draw[step=5mm, gridcol, very thin] (0,0) grid (\linewidth,#1);\end{tikzpicture}\par}
\newcommand{\aufgabe}[2]{\medskip\par\noindent{\color{accent}\large\textbf{Aufgabe #1}}\hfill{\color{accent2}\small #2}\par\vspace{0.3em}\hrule height0.4pt\vspace{0.6em}}
\newcommand{\kopf}[1]{\noindent{\color{accent}\Large\textbf{Numerik — #1}}\par\vspace{0.4em}\noindent\textbf{Name:}~\rule[-2pt]{5cm}{0.4pt}\hfill\textbf{Datum:}~\rule[-2pt]{3cm}{0.4pt}\par\vspace{0.3em}\hrule height0.8pt\vspace{1em}}
\setlength{\parindent}{0pt}
\begin{document}
\kopf{Numerische Integration}

\textbf{Notation der Vorlesung.} Eine Newton--Cotes-Formel approximiert
$\int_a^b f(x)\,dx\approx\sum_i w_i\,f(x_i)$ mit Stützstellen $x_i$ und Gewichten $w_i$.
Schrittweite $h=\frac{b-a}{n}$, Knoten $x_i=a+i\,h$. Wichtige Regeln (auf einem Intervall der Länge $H=b-a$):
\[
\text{Mittelpunkt: } H\!\cdot\! f\!\Big(\tfrac{a+b}{2}\Big),\quad
\text{Trapez: } \tfrac{H}{2}\big[f(a)+f(b)\big],\quad
\text{Simpson: } \tfrac{H}{6}\big[f(a)+4f(m)+f(b)\big].
\]
Monte-Carlo: $\hat I=|\Omega|\cdot\frac1N\sum_{i=1}^N f(X_i)$ mit $X_i\sim\mathrm{Uniform}(\Omega)$,
Standardfehler $\mathrm{SE}=\frac{|\Omega|\,\sigma_f}{\sqrt N}$.

\aufgabe{1}{leicht}
Gegeben sei das Integral $\displaystyle\int_{-2}^{-1}\sin x\,dx$ (exakter Wert
$\cos(-2)-\cos(-1)\approx-0{,}9564$).
\begin{enumerate}[label=(\alph*)]
  \item Approximieren Sie das Integral mit der \emph{Mittelpunktregel} bei \textbf{einem}
        Teilintervall: $I\approx (b-a)\,f\!\big(\tfrac{a+b}{2}\big)$.
  \item Verwenden Sie nun \textbf{zwei} gleich große Teilintervalle (zusammengesetzte
        Mittelpunktregel, $h=\tfrac12$). Werten Sie $\sin$ in den Mittelpunkten aus.
  \item Geben Sie für beide Fälle den absoluten Fehler zum exakten Wert an. Wie ändert sich
        der Fehler bei der Verfeinerung?
\end{enumerate}
\karo{7cm}

\aufgabe{2}{leicht}
Berechnen Sie $\displaystyle\int_0^1 (x+x^2)\,dx$ näherungsweise mit der \textbf{Trapezregel}
und Schrittweite $h=0{,}5$ (also Knoten $x_0=0,\ x_1=0{,}5,\ x_2=1$).
\begin{enumerate}[label=(\alph*)]
  \item Stellen Sie die zusammengesetzte Trapezformel
        $\frac h2\big[f(x_0)+2f(x_1)+f(x_2)\big]$ auf und rechnen Sie sie aus.
  \item Der exakte Wert ist $\tfrac56\approx0{,}8333$. Bestimmen Sie den Fehler.
  \item \emph{Skizzen-Hinweis:} Zeichnen Sie $f$ und die zwei Trapeze. Über- oder
        unterschätzt die Trapezregel hier das Integral? Begründen Sie mit der Krümmung von $f$.
\end{enumerate}
\karo{7cm}

\aufgabe{3}{mittel}
Die zusammengesetzte Trapezregel mit $n$ Teilintervallen lautet
\[
\int_a^b f\,dx\approx \frac h2\Big[f(x_0)+2\sum_{i=1}^{n-1}f(x_i)+f(x_n)\Big],\qquad h=\frac{b-a}{n}.
\]
\begin{enumerate}[label=(\alph*)]
  \item Erklären Sie, warum die \emph{inneren} Knoten $x_1,\dots,x_{n-1}$ das Gewicht $2$
        bekommen, die Randknoten $x_0,x_n$ aber nur das Gewicht $1$.
  \item Wie viele Funktionsauswertungen benötigt die Formel für $n$ Teilintervalle?
  \item Schreiben Sie die Formel für $n=4$ explizit aus (mit den konkreten Gewichten
        $1,2,2,2,1$).
\end{enumerate}
\karo{6cm}

\aufgabe{4}{mittel}
Approximieren Sie $\displaystyle\int_{-2}^{-1}\sin x\,dx$ mit der \textbf{Simpson-Regel (1/3)}
\[
\int_a^b f\,dx\approx \frac{b-a}{6}\big[f(a)+4f(m)+f(b)\big],\qquad m=\tfrac{a+b}{2}.
\]
\begin{enumerate}[label=(\alph*)]
  \item Berechnen Sie die Näherung von Hand (nutzen Sie $\sin(-2),\sin(-1{,}5),\sin(-1)$).
  \item Vergleichen Sie mit dem exakten Wert $-0{,}9564$. Wie groß ist der Fehler?
  \item Simpson ist exakt für Polynome bis Grad~$3$. Erklären Sie kurz, warum die Näherung
        hier (glatte Funktion) so genau ausfällt.
\end{enumerate}
\karo{7cm}

\aufgabe{5}{mittel}
Betrachten Sie $\displaystyle\int_0^1 e^{x}\,dx$ (exakt $e-1\approx1{,}71828$).
\begin{enumerate}[label=(\alph*)]
  \item Wenden Sie die \textbf{3/8-Regel} an: Knoten $0,\tfrac13,\tfrac23,1$, Gewichte
        $\tfrac18,\tfrac38,\tfrac38,\tfrac18$ (mal $(b-a)=1$).
  \item Wenden Sie die zusammengesetzte \textbf{Trapezregel} mit \emph{derselben} Schrittweite
        $h=\tfrac13$ an (Knoten $0,\tfrac13,\tfrac23,1$).
  \item Vergleichen Sie die Fehler. Warum ist dieser direkte Vergleich „nicht ganz fair''?
        (Hinweis: beide nutzen $4$ Stützstellen, aber die 3/8-Regel verteilt die Gewichte
        polynomial optimal.)
\end{enumerate}
\karo{7cm}

\aufgabe{6}{mittel}
Approximieren Sie $\displaystyle\int_0^1 e^{x}\,dx$ mit der \textbf{Boole-Regel}
(5 Stützstellen $0,\tfrac14,\tfrac12,\tfrac34,1$, Gewichte
$\tfrac{7}{90},\tfrac{32}{90},\tfrac{12}{90},\tfrac{32}{90},\tfrac{7}{90}$, mal $(b-a)=1$).
\begin{enumerate}[label=(\alph*)]
  \item Werten Sie $e^x$ an den 5~Knoten aus (auf 4 Nachkommastellen).
  \item Setzen Sie in die Boole-Formel ein und bestimmen Sie die Näherung.
  \item Bestimmen Sie den absoluten Fehler zum exakten Wert $e-1\approx1{,}71828$.
\end{enumerate}
\karo{7cm}

\aufgabe{7}{mittel--schwer}
Für $\displaystyle\int_0^1 e^{x}\,dx$ (exakt $1{,}71828$) sollen drei Verfahren auf
\emph{einem} Intervall $[0,1]$ verglichen werden: Trapez ($\tfrac12[f(0)+f(1)]$), Simpson
($\tfrac16[f(0)+4f(\tfrac12)+f(1)]$) und Boole (aus Aufgabe~6).
\begin{enumerate}[label=(\alph*)]
  \item Berechnen Sie alle drei Näherungen.
  \item Geben Sie für jedes Verfahren den \emph{prozentualen} Fehler
        $\frac{|\text{Näherung}-\text{exakt}|}{\text{exakt}}\cdot100\%$ an.
  \item Ordnen Sie die Verfahren nach Genauigkeit. Welcher Zusammenhang besteht zwischen
        Anzahl der Stützstellen und Genauigkeit?
\end{enumerate}
\karo{7cm}

\aufgabe{8}{mittel}
Bestimmen Sie das Interpolationspolynom $P_2(x)$ durch die drei Punkte
$(0,1),\ (0{,}5,3),\ (1,7)$ mit der \textbf{Lagrange-Interpolation}.
\begin{enumerate}[label=(\alph*)]
  \item Stellen Sie die Basispolynome $L_0(x),L_1(x),L_2(x)$ auf
        ($L_i(x)=\prod_{j\ne i}\frac{x-x_j}{x_i-x_j}$).
  \item Bilden Sie $P_2(x)=\sum_i y_i L_i(x)$ und fassen Sie zu einem Polynom
        $ax^2+bx+c$ zusammen.
  \item Prüfen Sie durch Einsetzen, dass $P_2$ tatsächlich durch alle drei Punkte geht.
\end{enumerate}
\karo{7cm}

\aufgabe{9}{schwer}
Newton-Cotes-Gewichte ergeben sich aus den Integralen der Lagrange-Basispolynome:
$w_i=\int_a^b L_i(x)\,dx$. Leiten Sie für die drei äquidistanten Knoten
$x_0=0,\ x_1=\tfrac12,\ x_2=1$ auf $[0,1]$ die Gewichte her.
\begin{enumerate}[label=(\alph*)]
  \item Schreiben Sie $L_0,L_1,L_2$ als Polynome aus.
  \item Berechnen Sie $w_i=\int_0^1 L_i(x)\,dx$ für $i=0,1,2$.
  \item Zeigen Sie, dass sich genau die Simpson-Gewichte $\tfrac16,\tfrac46,\tfrac16$ ergeben.
\end{enumerate}
\karo{8cm}

\aufgabe{10}{mittel--schwer}
\textbf{Monte-Carlo von Hand.} Das Integral $\displaystyle\int_1^3 x\,dx$ (exakt $4$) soll mit
$N=4$ Zufallsstichproben $X_i\sim\mathrm{Uniform}(1,3)$ geschätzt werden. Gegeben seien
\[
X_1=1{,}2,\quad X_2=1{,}8,\quad X_3=2{,}4,\quad X_4=2{,}9,\qquad f(x)=x.
\]
\begin{enumerate}[label=(\alph*)]
  \item Berechnen Sie den Schätzer $\hat I=(b-a)\,\frac1N\sum_i f(X_i)$ mit $|\Omega|=b-a=2$.
  \item Berechnen Sie die Stichprobenvarianz $\sigma_f^2=\frac{1}{N-1}\sum_i(f(X_i)-\bar f)^2$
        und den Standardfehler $\mathrm{SE}=\frac{(b-a)\,\sigma_f}{\sqrt N}$.
  \item Was geschieht mit dem Standardfehler, wenn man die Stichprobenzahl auf $4N$ vervierfacht?
\end{enumerate}
\karo{7cm}

\aufgabe{11}{schwer}
\textbf{Fluch der Dimensionen.}
\begin{enumerate}[label=(\alph*)]
  \item Eine Quadraturregel mit $N$ Punkten pro Achse braucht in $d$ Dimensionen $N^d$
        Auswertungen. Berechnen Sie die Punktzahl für $N=10$ und $d=1,3,10$.
  \item Monte-Carlo konvergiert mit $O(1/\sqrt N)$ \emph{unabhängig} von $d$. Um wie viel müssen
        Sie $N$ erhöhen, um den Fehler zu halbieren? Wie viele Punkte für eine Verzehnfachung
        der Genauigkeit?
  \item Erklären Sie die beiden Bezüge der Vorlesung: (i) die Mittelpunktregel ist ein
        Monte-Carlo-Schätzer mit $N=1$; (ii) der Mini-Batch-Gradient im SGD ist ein
        Monte-Carlo-Schätzer des vollen Gradienten.
\end{enumerate}
\karo{8cm}

\end{document}
